\section{Trusted Data: gobernanza de datos para modelos de alto riesgo}

La sección anterior dependía de etiquetas de entrenamiento y de verified outcome; esta sección indaga de dónde provienen esos datos, quién puede acceder a ellos y cómo demostrar que son legales y confiables. Los datos de seguridad de alto riesgo son a la vez escasos y sensibles, por lo que no pueden seguir el método de la extracción común de internet ni el de «primero copiar al data lake y después gobernar». La barrera central de la capacidad de datos no es la cantidad de archivos, sino la autorización, la provenance, la segregation, el label protocol y la auditoría continua.

\subsection{Provenance, segregation y retention}

La \term{provenance} (origen y linaje de procesamiento) registra quién aportó los datos, bajo qué autorización y mediante qué transformaciones llegaron a una versión de entrenamiento o de evaluación. La \term{data segregation} (aislamiento de datos) separa los datos de sensibilidad extrema de los registros generales, los entornos de desarrollo, las cuentas de bajos privilegios y otros inquilinos. La \term{retention} (política de conservación) establece por qué se conservan los datos, durante cuánto tiempo y cuándo se eliminan o pasan a legal hold, en lugar de conservarlos de manera permanente por omisión.

\lecturefigure{06-trusted-data.png}{Los datos de entrenamiento de alto riesgo requieren fuentes confiables, procesamiento trazable, acceso aislado y una política de conservación explícita.}{Subtítulos locales 16:20--20:30; contexto de la colaboración entre Thorn y NCMEC; redibujo conceptual.}

\begin{knowledgebox}{Lectura de la figura: la gobernanza de datos abarca todo el lifecycle}
La trusted source demuestra la autorización y el uso; la provenance vincula la evidencia original, las transformaciones de desidentificación, la label y la dataset version; la segregation restringe entornos, redes y personal; la retention/audit se ocupa de la eliminación, la respuesta a incidentes y las inspecciones regulatorias. Si falta cualquiera de estas piezas, el equipo podría no ser capaz de explicar por qué un modelo aprendió cierto patrón, por qué un empleado pudo acceder a los datos, o si una solicitud de eliminación realmente alcanzó a los derived artifacts.
\end{knowledgebox}

\begin{importantbox}{Label quality no equivale a votación por mayoría}
Las etiquetas de alto riesgo requieren límites claros fijados por definiciones legales, la policy de la plataforma y las guías operativas. El conjunto de entrenamiento debe distinguir entre verified positive, policy violation, uncertain, benign hard negative y unavailable context; la evaluación además debe registrar el reviewer agreement y la adjudication. Mezclar todo el contenido «incómodo» en una sola clase positiva produce modelos imposibles de interpretar y amplía los daños colaterales contra la expresión legítima.
\end{importantbox}

\teachervoice{El ponente enfatiza que Thorn puede entrenar este tipo de modelos no porque haya encontrado un conjunto de datos público más grande, sino porque, mediante una colaboración confiable de largo plazo, obtuvo material controlado, legal y con contexto. Nota de clase: en dominios sensibles, el data access es en sí mismo una responsabilidad de gobernanza, no una ventaja común de adquisición.}

\subsection{Los datos de entrenamiento, evaluación y producción deben tener funciones separadas}

Esta sección separa los usos de los datos para evitar que el equipo del modelo ajuste parámetros una y otra vez sobre el mismo reference set y aun así presente el resultado como una prueba independiente. Los datos de entrenamiento sirven para ajustar los parámetros; los datos de validation, para elegir umbrales y modelos; los datos de test, para la evaluación final; y el production feedback solo puede pasar a la siguiente versión después de analizar sus sesgos y sus mecanismos de selección. Para cada conjunto deben fijarse el snapshot, la versión, los roles de acceso y las verificaciones de contamination.

\begin{warningbox}{El verified feedback también puede tener selection bias}
Solo los objetos que el modelo envía a revisión tienden a recibir etiquetas, por lo que los datos de retroalimentación se inclinan de manera natural hacia las regiones que el modelo actual considera «sospechosas». Si se reentrena directamente con esas muestras, el sistema podría volverse cada vez más ciego a los riesgos que quedan fuera de los umbrales anteriores. El equipo necesita muestreo aleatorio, shadow evaluation, un reference set independiente y red-team evidence externa para descubrir los blind spots.
\end{warningbox}

\subsection{Resumen de la sección}

Los trusted data son la combinación de autorización legal, aislamiento técnico, protocolos de etiquetado y evidencia de auditoría. Si un modelo de alto riesgo no puede explicar el origen, el uso y la ruta de eliminación de sus datos de entrenamiento, no cumple las condiciones para un despliegue seguro, aunque sus métricas fuera de línea sean muy altas.
